% !TeX root = ../drafts/02-vm-specification.tex
\section{VM specification}\label{sec:vm}

\paragraph{Field.}
\[
  \K=\F_{2^{64}}=\F_2[x]/(x^{64}+x^4+x^3+x+1),
  \qquad
  \E=\K[y]/(y^3+y+1),
  \qquad |\E|=2^{192}.
\]
A \emph{memory word}, including a cell value, immediate, or hash value, is an element $v=v_0+v_1y+v_2y^2\in\E$ with $v_0,v_1,v_2\in\K$. Each memory address, program counter $\pc$, frame pointer $\fp$, and read counter is in $\K$.

\paragraph{Addresses as powers of a generator.} Fix a generator $\gen$ of the multiplicative group $\K^{\times}$, of order $2^{64}-1$. Memory addresses and bytecode addresses (the program counter) are powers of $\gen$: the $i$-th address is the $\K$-element $\gen^{\,i}$.

\paragraph{Memory.} Memory is read-only (equivalently: write-once): an array of $2^{h}$ cells, each holding one $192$-bit word (an $\E$-element), addressed by the $2^{h}$ powers $\gen^{0},\gen^{1},\dots,\gen^{2^{h}-1}$ (the addresses are $\K$-elements). The memory log-size $h$ depends on each VM execution, with the constraint $16\le h\le 32$.

\paragraph{Program.} The program is a fixed, public sequence of instructions, addressed like memory by powers of $\gen$: the program counter $\gen^{\,i}$ selects the $i$-th instruction. An instruction is an opcode together with its operands, described below.

\paragraph{Registers.} The machine state is two registers, the program counter $\pc$ and the frame pointer $\fp$, each a power of $\gen$. Every instruction except \texttt{JUMP} advances the state by $\pc\gets\gen\cdot\pc$ (fetching the next instruction) and leaves $\fp$ unchanged; \texttt{JUMP} may set both.

\paragraph{Operands and addressing.} An operand is either an \emph{immediate}, a literal $\E$-element used as given, or an \emph{fp-relative reference}: an offset $j$ into the current frame, encoded as the $\gen$-power $o=\gen^{\,j}$. With the frame based at $\fp=\gen^{\mathrm{base}}$, the reference names the cell
\[
  \mem[\fp\cdot o],\qquad \fp\cdot o\;=\;\gen^{\mathrm{base}}\cdot\gen^{\,j}\;=\;\gen^{\,\mathrm{base}+j}.
\]

\paragraph{Instruction set.} The machine implements six instructions. Below, $o,o_A,\dots$ are fp-relative references, $k$ is an immediate operand, and $\loc{o}$ is the cell at address $\fp\cdot o$.
\begin{center}
\begin{tabularx}{\textwidth}{@{}llX@{}}
\hline
Instruction & Operands & Semantics\\
\hline
\texttt{XOR}           & $[o_A,o_B,o_C]$ & $\loc{o_C}=\loc{o_A}+\loc{o_B}$ \quad(addition in $\E$)\\
\texttt{MUL\_NATIVE}   & $[o_A,o_B,o_C]$ & $\loc{o_C}=\loc{o_A}\cdot\loc{o_B}$ (multiplication in $\E$)\\
\texttt{SET\_CONSTANT} & $[o,k]$         & $\loc{o}=k$ \quad($k$ a $192$-bit immediate)\\
\texttt{DEREF}         & $[\alpha,\beta,\gamma;\,s]$ & $\mem[\loc{\alpha}\cdot\beta]=\srcsel{s}$\\
\texttt{JUMP}          & $[o_c,o_d,o_f]$ & conditional jump (below)\\
\texttt{SHA2}          & $[o_0,o_1,o_2,o_3,o_V,o_C]$ & SHA-256 compression (below)\\
\hline
\end{tabularx}
\end{center}

\texttt{SHA2} is the SHA-256 compression $C$ \cite{fips180-4}, consuming a $64$-byte message block and a $32$-byte chaining value and producing $32$ bytes. Each $128$-bit chunk occupies one full $\E$ memory cell through the canonical embedding $a_0+a_1y\mapsto a_0+a_1y+0y^2$; using a cell as a \texttt{SHA2} operand therefore constrains its top limb to zero. The four message chunks are addressed \emph{independently} by references $o_0,\dots,o_3$ (chunk $i$ at $\loc{o_i}=\fp\cdot o_i$), while $o_V$ names two \emph{consecutive} chaining-value cells $\loc{o_V},\loc{\gen\cdot o_V}$. The $256$-bit result is written to the consecutive cells $\loc{o_C},\loc{\gen\cdot o_C}$, so each row accesses eight cells. The instruction carries no immediate.

The hash the machine actually computes is \shaeth, a length-prefixed Merkle--Damg\aa{}rd over $C$. For a message of $n$ bits,
\[
  \shaeth(m)=C\bigl(\cdots C(C(\ivbase,\lenblock(n)),m_1)\cdots,m_b\bigr),
\]
where $\ivbase=\mathrm{SHA256}(\texttt{"Ethereum"})$, $\lenblock(n)$ is $n$ as a $512$-bit big-endian integer, and $m_1,\dots,m_b$ are the $64$-byte blocks of $m$ zero-padded up to a whole number of blocks. Prefixing the length rather than appending it makes the encoding prefix-free, hence the construction indifferentiable from a random oracle; the usual $10^\ast$ padding is then unnecessary, the length block already separating messages of different lengths and the zero fill being injective at a fixed one.

Prefixing is also what makes the machine's hash free of overhead. The first compression absorbs nothing but the length, so
\[
  \ivlen(n)=C(\ivbase,\lenblock(8n))
\]
depends on $n$ alone. Every call site in a program knows its length at assembly time, so that compression is a constant the bytecode supplies through $o_V$, never an executed instruction: an $n$-byte hash is exactly $\lceil n/64\rceil$ \texttt{SHA2} rows. A one-block $64$-byte hash, which is what every sponge step and Merkle node is, therefore takes the chaining value $\ivlen(64)$ and one row; a longer one passes each result back through $o_V$. There is no byte counter, no finalization flag, and no chunk tree: this is why a single compression opcode is a complete hash primitive here. The \texttt{SHA2} table (\S\ref{sec:tab-sha2}) wires its memory and bytecode operands into a flock~\cite{flock-impl} R1CS proof.

\texttt{DEREF} stores, at the dereferenced address $\loc{\alpha}\cdot\beta$, a value chosen by a \emph{store mode} $s\in\{\texttt{deref\_cell}[\gamma],\texttt{deref\_pc},\texttt{deref\_fp}\}$:
\[
  \srcsel{s}\;=\;
  \begin{cases}
    \loc{\gamma} & s=\texttt{deref\_cell}[\gamma] \quad(\text{the local cell at offset }\gamma),\\
    \gen^{2}\cdot\pc & s=\texttt{deref\_pc} \quad(\text{a return address: the program counter advanced by two}),\\
    \fp & s=\texttt{deref\_fp} \quad(\text{the current frame pointer}).
  \end{cases}
\]
The \texttt{deref\_pc} and \texttt{deref\_fp} modes let a caller save its return address and frame so a callee can later return and restore them, the basis for function calls. The skip of two is the call's return target: place the \texttt{DEREF} that saves $\gen^{2}\cdot\pc$ immediately before the \texttt{JUMP} that transfers control, and $\gen^{2}\cdot\pc$ names the instruction just after that \texttt{JUMP}.

\texttt{JUMP} reads a condition $c=\loc{o_c}$, a $\K$-element, and branches on whether it is zero. When $c\neq0$ it transfers control, $\pc\gets\loc{o_d}$ and $\fp\gets\loc{o_f}$; when $c=0$ it falls through, $\pc\gets\gen\cdot\pc$ and $\fp\gets\fp$.

\paragraph{Words used as addresses.} A memory word is $192$-bit, while $\pc$, $\fp$, and addresses are in $\K$. Some instructions (\texttt{DEREF}, \texttt{JUMP}) interpret memory words as address or registers. The VM then requires those words to be stricly in $\K$ (i.e. the two top lanes are zero). For \texttt{JUMP} the requirement is unconditional: $\loc{o_d}$ and $\loc{o_f}$ must be in $\K$ on every execution of the instruction, not only when the branch is taken. Its condition is $\K$-valued too, though nothing reads it as an address; all three are committed one lane each, the memory interaction zeroing the rest (\S\ref{sec:tab-jump}).
